% SPDX-License-Identifier: CC-BY-4.0
% Copyright (c) 2026 Martin Schröder <info@swedishembedded.com>

\section{System}\label{sec:system}

The loop is a composition of three existing projects with one rule of ownership
(\cref{fig:arch}): Sven runs agents and records what they did; Brain computes
(inference, LoRA training, serving); Splinter decides what to learn, what counts as
evidence, what to train on, and what to release. Splinter does not re-implement either
dependency.

\begin{figure}[t]
\centering
\begin{tikzpicture}[
  box/.style={draw,rounded corners,align=center,minimum width=3.3cm,minimum height=1.0cm,font=\small},
  arr/.style={-{Stealth[length=2mm]},thick}]
\node[box,fill=blue!8] (src) {Repository\\(git revision)};
\node[box,fill=green!10,right=1.1cm of src] (gen) {Task generation\\+ validation};
\node[box,fill=green!10,right=1.1cm of gen] (solve) {Solve in a\\workspace};
\node[box,fill=orange!15,below=1.0cm of solve] (sven) {Sven\\(agent, tools, ATIF)};
\node[box,fill=green!10,below=1.0cm of gen] (ver) {Hidden-test\\verifier + exposure audit};
\node[box,fill=green!10,below=1.0cm of src] (store) {Experience store\\(immutable, lineage)};
\node[box,fill=green!10,below=1.0cm of store] (data) {Dataset\\projection};
\node[box,fill=red!10,right=1.1cm of data] (train) {Brain\\LoRA training};
\node[box,fill=green!10,right=1.1cm of train] (rel) {Release gate\\+ manifest};
\draw[arr] (src) -- (gen);
\draw[arr] (gen) -- (solve);
\draw[arr] (solve) -- (sven);
\draw[arr] (sven) -- (ver);
\draw[arr] (ver) -- (store);
\draw[arr] (store) -- (data);
\draw[arr] (data) -- (train);
\draw[arr] (train) -- (rel);
\draw[arr,dashed] (rel.north) |- ([yshift=0.5cm]solve.north) -- (solve.north);
\end{tikzpicture}
\caption{The loop. Green: Splinter; orange: Sven; red: Brain. The dashed arrow is the
next round: rollouts run on the released adapter. Hidden tests exist only in the
verifier's own directory after an episode ends.}\label{fig:arch}
\end{figure}

\paragraph{Workspace tasks.} A task is a snapshot of a repository tree at a revision, an
instruction stating a symptom, and privileged material the solver never receives:
the hidden checks (fail-to-pass and pass-to-pass commands) and the paths that carry
them. The environment records the snapshot digest and the sandbox, so two records of
the same snapshot and sandbox have the same environment digest. This is a new
\emph{environment kind} and a new \emph{verifier kind} of an existing vocabulary, not
a software-specific schema; document, tool and person tasks use the same experience
store, curriculum, release gate and lineage.

\paragraph{Validity before use.} A candidate task is admitted only if it passes the
sequence of \cref{fig:valid}: the clean snapshot's selected tests pass twice (flaky
baselines are refused), the mutation applies, at least one previously passing test now
fails reproducibly, and the reverse patch restores green. With
$T_{\mathrm{base}}$ the tests passing before the change and $T_{\mathrm{bug}}$ those
passing after, $\mathrm{F2P}=T_{\mathrm{base}}\setminus T_{\mathrm{bug}}$ must be
non-empty. The failing tests are moved out of the snapshot the solver receives. A
repair passes when $\mathrm{F2P}$ passes with the hidden files overlaid fresh and the
pass-to-pass tests still pass; a patch that touches a hidden or protected path fails.

\begin{figure}[t]
\centering
\begin{tikzpicture}[
  stp/.style={draw,rounded corners,align=center,minimum width=2.2cm,minimum height=0.9cm,font=\scriptsize},
  arr/.style={-{Stealth[length=1.6mm]},thick},
  rej/.style={draw,dashed,rounded corners,align=center,font=\scriptsize,text=red!60!black}]
\node[stp] (a) {Baseline\\passes $\times2$};
\node[stp,right=0.55cm of a] (b) {Apply\\change};
\node[stp,right=0.55cm of b] (c) {Some test\\fails (F2P)};
\node[stp,right=0.55cm of c] (d) {Failure\\reproducible};
\node[stp,right=0.55cm of d] (e) {Reverse patch\\restores green};
\node[stp,right=0.55cm of e,fill=green!12] (f) {Seal hidden\\files};
\foreach \x/\y in {a/b,b/c,c/d,d/e,e/f} \draw[arr] (\x) -- (\y);
\node[rej,below=0.5cm of a] (ra) {flaky or red:\\refuse};
\node[rej,below=0.5cm of c] (rc) {no failing test\\(equivalent): refuse};
\node[rej,below=0.5cm of e] (re) {not restorable:\\refuse};
\draw[arr,dashed,red!60!black] (a) -- (ra);
\draw[arr,dashed,red!60!black] (c) -- (rc);
\draw[arr,dashed,red!60!black] (e) -- (re);
\end{tikzpicture}
\caption{Admission of a mutation-seeded task. Every refusal is counted by reason.}\label{fig:valid}
\end{figure}

\paragraph{Solving and recording.} The solver is a Sven coding worker with file,
search and shell tools, a tool-call budget, a repetition stop and a deadline, started in a
fresh copy of the snapshot; visible checks may give it follow-up turns. The recorded
unit is Sven's full trajectory in ATIF: assistant turns, tool calls, observations. The
verdict is produced outside the agent. The tool set, system prompt, budgets and
sampling are one pinned \emph{harness profile} used identically for the base and the
adapter.

\paragraph{Training data.} Only trajectories whose patch passes the verifier and whose
exposure audit is clean are eligible. Each assistant decision becomes a training record
conditioned on the prior conversation with the tool definitions the agent was offered;
user, tool and system turns are context with no loss. Trajectories with compaction, with
unparsed tool arguments or over the training block length are excluded before training,
with counted reasons.

\paragraph{Order that makes the exam honest.} Families are assigned to training or exam
\emph{before} any rollout (\cref{fig:split}); exam families are never solved for data
purposes, and a family is a mutated path or source entity, so near-copies of one defect
stay on one side.

\begin{figure}[t]
\centering
\begin{tikzpicture}[
  box/.style={draw,rounded corners,align=center,minimum height=0.8cm,font=\small},
  arr/.style={-{Stealth[length=2mm]},thick}]
\node[box,fill=blue!8,minimum width=2.6cm] (all) {Admitted tasks\\grouped by family};
\node[box,fill=gray!15,right=1.4cm of all,yshift=0.9cm,minimum width=3.0cm] (ex) {Exam families\\(reserved first)};
\node[box,fill=green!10,right=1.4cm of all,yshift=-0.9cm,minimum width=3.0cm] (tr) {Training families};
\node[box,fill=green!10,right=1.0cm of tr,minimum width=2.6cm] (ro) {Rollouts $\to$ verify\\$\to$ dataset $\to$ LoRA};
\node[box,fill=orange!15,right=1.0cm of ex,minimum width=2.6cm] (ev) {Base and adapter\\same seeds, same harness};
\draw[arr] (all) -- (ex);
\draw[arr] (all) -- (tr);
\draw[arr] (tr) -- (ro);
\draw[arr] (ex) -- (ev);
\draw[arr,dashed] (ro) |- ([yshift=-0.6cm]ev.south) -- (ev.south);
\end{tikzpicture}
\caption{Split before rollout. Nothing from an exam family enters a rollout used for data.}\label{fig:split}
\end{figure}

\paragraph{Release.} The release is an immutable adapter with a manifest naming its
base checkpoint by digest; the served alias moves by compare-and-set and can be rolled
back. Sven then talks to the fine-tuned model through Brain's OpenAI-compatible server
with that manifest.

\begin{table}[t]
\centering\small
\begin{tabular}{p{4.6cm}p{2.4cm}p{6.2cm}}
\toprule
Component & State & Note \\
\midrule
Coding worker, trajectory read-back & present & used by an existing coding-loop sample \\
Per-action training view, LoRA, release gate & present & gate probes are closed-book only (\cref{sec:audit}) \\
Workspace environment and sandbox runs & in progress & tree in, command run, diff out \\
Tool-turn loss masking in the trainer & in progress & see \cref{sec:audit} \\
Trajectory fidelity in the agent SDK & in progress & compaction, offered tools, raw arguments \\
Adapter manifest writer & in progress & needed for serving a release \\
Workspace verifier, exposure audit & not started & \\
Task generator (mutation-seeded, import) & not started & \\
Agentic release probes, serve check & not started & \\
\bottomrule
\end{tabular}
\caption{Implementation state at the date of this version.}\label{tab:state}
\end{table}
